% Original starter content for the Sedimentologika class.
\documentclass[preprint]{sedimentologika}
\englishvariety{british}
\preprintdeclaration{This is a fictional pre-print starter; replace this declaration before submission.}

\title{A transparent workflow for sediment observations}
\shorttitle{Transparent sediment observations}
\author[1]{Maya Rivera}
\author[1]{Jon Bell}
\authorshdr{M. Rivera and J. Bell}
\affil[1]{Northbridge Earth Systems Lab, Example University}
\correspondingauthor{Maya Rivera}{maya.rivera@example.edu}
\CRediT{Writing-original draft}{MR}
\CRediT{Writing-review \& editing}{JB}
\CRediT{Conceptualization}{MR, JB}
\CRediT{Data curation}{JB}
\CRediT{Formal analysis}{MR}
\CRediT{Investigation}{MR, JB}
\CRediT{Methodology}{MR, JB}
\CRediT{Visualization}{MR}
\keywords{sedimentology\\field methods\\provenance\\open data}
\addbibresource{references.bib}

\begin{document}
\SedKtitle{
  This study demonstrates how a field team can keep sediment observations,
  processing decisions, and data provenance together in one readable record.
}{
  We present a small workflow for classifying sediment beds from repeatable
  field observations.  The starter is intentionally fictional and is meant to
  be replaced with a real study while preserving the class's required
  declarations.
}{}

\section{Introduction}
Sediment descriptions are most useful when observations can be traced from a
field notebook to the reported interpretation.  A short provenance ledger
reduces ambiguity about which observations were retained and why.

\section{Methods}
We record grain-size class, sorting, bedding structure, and confidence at each
station.  Each station receives a stable identifier before interpretation.
The workflow is designed to be explicit enough for a collaborator to repeat,
and the data and code statement below shows where the real project links would
be placed.

\begin{table}[ht]
  \centering
  \caption{Example station register.}
  \label{tab:stations}
  \begin{tabular}{lccc}
    \toprule
    Station & Median class & Sorting score & Confidence \\
    \midrule
    S01 & fine sand & 0.82 & high \\
    S02 & medium sand & 0.64 & medium \\
    S03 & granule & 0.47 & medium \\
    \bottomrule
  \end{tabular}
\end{table}

\section{Results}
The station register makes the transition from observation to interpretation
visible.  The largest uncertainty is associated with mixed grain populations,
so those observations should be retained as mixtures rather than forced into
a single class.  The table is a compact example of the kind of data summary
that can be expanded with figures and supplementary material.

\section{Discussion}
Open workflows benefit from recording both positive and negative decisions:
which observations were excluded, which measurements were recalibrated, and
which interpretations remain provisional.  A structured manuscript can carry
those decisions alongside the scientific narrative.

\section{Conclusions}
The `sedimentologika` class provides a journal-oriented preprint layout with
plain-language summary, CRediT roles, declarations, and bibliography support.
Replace the fictional station values and links before submission.

\acknowledgements{Thanks to the field team and to the maintainers of the open tools used in this example.}
\datacode{The fictional observations and code for this starter are not an empirical dataset. Replace this statement with a DOI or repository link for the real study.}
\supplementarymaterial{No supplementary material accompanies this fictional starter.}
\conflictofinterest{The authors declare no conflicts of interest.}
\artificialintelligence{No generative AI tools were used to collect or analyse the fictional observations.}
\funding{No external funding supported this fictional example.}
\printbibliography
\end{document}
